
Soit @@M0@@ l'homomorphisme ainsi défini. Si l'on désigne par @@M1@@ l'idéal @@M2@@ de @@M3@@, on a @@M4@@. Gardant les notations du lemme 2.1.4 on trouve un homomorphisme de groupes abéliens
@@M5@@
qui associe à la classe de tout @@M6@@-module inversible @@M7@@ dans @@M8@@ l'élément @@M9@@ de @@M10@@.

On tire immédiatement des lemmes 2.1.2 et 2.1.4 :

\textbf{Proposition 2.2.1.} \uline{Pour tout schéma noethérien} @@M11@@ \uline{et pour tout} @@M12@@ (@@M13@@\uline{) l'application composée}
@@M14@@
\uline{définie par}
@@M15@@
\uline{est une application polynôme de degré} @@M16@@.

\subsection*{2.3.}
Si @@M17@@ est un schéma propre sur un corps @@M18@@ et si @@M19@@ est le morphisme structural, alors l'homomorphisme
@@M20@@
associe à la classe @@M21@@ de tout @@M22@@-module cohérent @@M23@@ la caractéristique d'Euler-Poincaré @@M24@@ de @@M25@@ sur @@M26@@. De 2.2.1 on conclut immédiatement :

